\section{Introduction}\label{sec:introduction}
Vision-language-action (VLA) models provide useful visual and language representations for robot control. Transferring these representations to a new robot still requires learning how its actions affect the scene. If the initial policy rarely completes the task, binary success provides little guidance for reinforcement learning.

Consider \emph{put the apple on the plate}. A failed rollout may contain a correct approach and grasp followed by a poor release. Treating the entire attempt as a failure discards this distinction. Training the parts separately helps exploration, but a policy can forget earlier skills while learning later ones. It must then connect the retained skills into a complete behavior.

We address these problems through the pipeline in Fig.~\ref{fig:hero}. Dense rewards first guide skill acquisition. Binary refinement then targets completion, while balanced self-imitation retains successful behavior across instructions. Finally, milestone-based RL trains one policy to execute the full task from an empty-gripper start.

The pipeline separates feedback design and skill consolidation from the robot-specific action mapping. We implement it with SmolVLA~\cite{smolvla}, LoRA, and a residual action controller. Experiments examine complete-task learning and skill retention in cable-driven parallel robot (CDPR) simulation, and intermediate feedback with OpenVLA-OFT~\cite{openvlaoft} on LIBERO-Spatial. An ongoing WX200 study extends the evaluation to a 5-DoF serial manipulator.

We contribute a training schedule that combines dense skill acquisition, binary refinement with instruction-balanced self-imitation, and full-task learning under one instruction. We specify milestone credit that retains local credit for achieved approach and pickup stages when later placement fails. Our experiments quantify skill forgetting and partial recovery, report complete-task performance after action-interface adaptation, and compare intermediate and terminal feedback on LIBERO-Spatial.
